% ====================================================================
%  CHAPTER 2: LITERATURE REVIEW
%  Two levels of subheading are usually enough. Each section should end
%  with what remains unresolved, so that the knowledge gaps in the final
%  section follow from the review rather than being asserted.
% ====================================================================
\chapter{LITERATURE REVIEW}\label{ch:2}

\section{[The Subject and Its Importance]}\label{sec:2.1}

\subsection{[Global and Regional Context]}\label{sec:2.1.1}
{[}Synthesise, do not list. Group studies by what they establish, and say where they disagree \parencite{example_2020_article, example_2021_chapter}.]

\subsection{[Local Context]}\label{sec:2.1.2}
{[}Evidence specific to your country, crop, system or population.]

\section{[The Constraint Your Study Addresses]}\label{sec:2.2}

\subsection{[First Dimension of the Constraint]}\label{sec:2.2.1}
{[}Text.]

\subsection{[Second Dimension of the Constraint]}\label{sec:2.2.2}
{[}Text.]

\section{[Approaches That Have Been Tried]}\label{sec:2.3}
\guidance{Assess each approach on its own terms: what it achieved, under what conditions, and why it has not resolved the problem. Distinguish evidence from extrapolation.}

{[}Text, with a comparative table if it helps. See Table~\ref{tab:2.1}.]

\begin{table}[htbp]
  \centering
  \caption{[Comparison of approaches reported in the literature].}
  \label{tab:2.1}
  \footnotesize
  \begin{tabularx}{\textwidth}{@{}p{0.22\textwidth} p{0.20\textwidth} X@{}}
    \toprule
    \textbf{Approach} & \textbf{Reported outcome} & \textbf{Limitation} \\
    \midrule
    {[}Approach 1] & {[}Outcome] & {[}Limitation] \\
    {[}Approach 2] & {[}Outcome] & {[}Limitation] \\
    \bottomrule
  \end{tabularx}
\end{table}

\section{Methodological Basis for the Study}\label{sec:2.4}
\guidance{Justify the methods you will use in Chapter 3 against alternatives. This is where you show that the design is considered rather than conventional.}

{[}Text.]

\section{Knowledge Gaps Addressed by the Study}\label{sec:2.5}
\guidance{State each gap as a numbered or clearly separated claim, and make sure each one maps onto a specific objective. Close with what the study will not answer.}

{[}First gap.] [Second gap.] [Third gap.]

{[}The study addresses these gaps. It will not answer \ldots]
